\documentclass[11pt]{article}
\usepackage[T1]{fontenc}
\usepackage{textcomp}
\usepackage[margin=0.75in]{geometry}
\usepackage{amsmath,amssymb,amsthm}
\usepackage{array,booktabs}
\usepackage{hyperref}
\setlength{\parindent}{0pt}
\newtheorem{theorem}{Theorem}
\theoremstyle{definition}
\newtheorem{definition}{Definition}
\title{Flow Proof Artifact: prop-07-square-whole-plus-segment}
\author{Generated by \texttt{flow doc proof}}
\date{Flow Proof Book}
\begin{document}
\maketitle
If a straight line is cut at random, the square on the whole and that on one segment equal twice the rectangle by the whole and that segment together with the square on the remaining segment.\\[0.5em]
\textbf{Source.} euclid — Elements, Book II, Proposition 7\\[0.5em]
\bigskip
\noindent{\large \textbf{Derived fact 1.} \textit{If a straight line is cut at random, the square on the whole and that on one segment equal twice the rectangle by the whole and that segment together with the square on the remaining segment} \hfill the Euclidean plane \cdot Euclid Book II \cdot Proposition 7: the square on the whole plus that on one segment equals twice a rectangle plus a square}\par
\smallskip\noindent\textit{Coordinate: the Euclidean plane \cdot Euclid Book II \cdot Proposition 7: the square on the whole plus that on one segment equals twice a rectangle plus a square}\par
\smallskip\noindent\textit{Source: euclid — Elements, Book II, Proposition 7}\par
\smallskip\noindent\textit{Built on: proposition 4: the square on the whole equals the squares on the parts plus twice their rectangle, for Euclid Book II on the Euclidean plane}\par
\medskip
\noindent\textbf{Goal.} If a straight line is cut at random, the square on the whole and that on one segment equal twice the rectangle by the whole and that segment together with the square on the remaining segment\par
\noindent$\displaystyle sq(whole) + sq(C) = 2 \cdot rect(whole, C) + sq(D)$\par
\medskip
\noindent
\begin{tabular}{@{} >{\raggedright\arraybackslash}p{0.06\textwidth} >{\raggedright\arraybackslash}p{0.47\textwidth} >{\raggedleft\arraybackslash}p{0.41\textwidth} @{}}
\textbf{\#} & \textbf{Proof} & \textbf{Mathematics} \\
\midrule
\textbf{1.} & We prove that proposition 7: the square on the whole plus that on one segment equals twice a rectangle plus a square for Euclid Book II the Euclidean plane. & \\[0.45em]
\textbf{2.} & Let whole = straight\_line\_AB. & \\[0.45em]
\textbf{3.} & Let C = segment\_AC. & \\[0.45em]
\textbf{4.} & Let D = segment\_CB. & \\[0.45em]
\textbf{5.} & From step 2, step 3, and step 4, this implies sq(whole) plus sq(C) equals 2 times rect(whole, C) plus sq(D). Hence proven. & $\displaystyle sq(whole) + sq(C) = 2 \cdot rect(whole, C) + sq(D)$ \\[0.45em]
\textbf{6.} & We invoke the derived fact governing Euclid Book II the Euclidean plane: proposition 4: the square on the whole equals the squares on the parts plus twice their rectangle, for Euclid Book II on the Euclidean plane. & \\[0.45em]
\bottomrule
\end{tabular}
\label{thm:Geometry-Euclid Book II-proposition_7_the_square_on_the_whole_plus_that_on_one_segment_equals_twice_a_rectangle_plus_a_square}
\par\medskip\hrule
\end{document}
